\section{PD-PPO Scheduler and Evaluation Protocol}
\label{sec:framework_protocol}

\subsection{Scheduler inputs, actions, and constraints}

The PD-PPO policy is trained using the policy reward in
\Cref{eq:training_reward}, whose forecast-loss term is derived from the
multi-step forecast loss in \Cref{eq:step_forecast_loss}. The macro-averaged
normalized forecast loss in \Cref{eq:macro_loss} is used for
event-type-balanced evaluation and does not enter PPO training. At each time step,
the scheduler receives recent sample-and-hold values and observation masks,
together with the time since each channel's last sampling opportunity. Its input
also includes channel operating states, the previous executed activation vector,
budget and time features, and simulated warning signals.
The reported actor and critic do not receive the remaining activation time
$r_t$ explicitly; the environment enforces the lock after an action is proposed.

The finite candidate action set is listed in
\Cref{tab:action_space_instantiation}. A feasibility check removes actions that
violate either the steady-state or startup-peak power limit and thereby forms the
feasible-action mask. During training, the policy samples an index from the
feasible action-index set. During evaluation, it selects the highest-scoring
feasible index. The resulting index specifies the proposed activation vector.
The minimum dwell-time constraint in \Cref{eq:minimum_activation_rule} is then
applied to determine the executed activation vector. Unless stated otherwise,
all baselines use the same feasibility check and execution rule.

\Cref{tab:action_space_instantiation} gives the candidate channel subsets used
by PD-PPO and the constrained baselines. Every subset includes the meteorological core
channel, and each differs only in whether and which optional channel is active. The evaluated
configuration therefore keeps one meteorological core channel active and allows at
most one optional channel to operate.

\begin{table}[htbp]
  \centering
  \caption{Candidate action set with a mandatory meteorological core channel and an
  optional-channel capacity of one.}
  \label{tab:action_space_instantiation}
  \footnotesize
  \setlength{\tabcolsep}{3.0pt}
  \begin{tabular}{@{}c p{0.35\linewidth} c p{0.21\linewidth} c@{}}
    \toprule
    Action index & Candidate channel subset & Power & Interpretation & \shortstack{Fixed-baseline\\selections} \\
    \midrule
    0 & Meteorological core channel & 0.25 & Core channel only & 0 of 22 \\
    1 & Meteorological core + basic radiometer channels & 0.33 & Optional channel & 0 of 22 \\
    2 & Meteorological core + shielded temperature--humidity channels & 0.43 & Optional channel & 0 of 22 \\
    3 & Meteorological core + infrared snow-surface-temperature channels & 0.74 & Thermal specialist channel & 9 of 22 \\
    4 & Meteorological core + laser disdrometer channels & 0.74 & Particle specialist channel & 2 of 22 \\
    5 & Meteorological core + FC4 snow mass-flux channels & 0.74 & Flux specialist channel & 11 of 22 \\
    \bottomrule
  \end{tabular}
  \vspace{0.4ex}

  \noindent\footnotesize
  Every candidate channel subset includes the meteorological core channel. The final
  column reports fixed baselines selected on the calibration/validation partition for the 22
  post-selection evaluation seeds. The scheduling decision selects at most one
  optional channel; action 0 selects none.
\end{table}


In the reported $q=1$ instance, the largest candidate steady-state demand is
$0.74<B=0.75$, and the largest startup-peak demand is $0.92<B_{\mathrm{peak}}=0.95$.
Consequently, every one of the six candidate action indices belongs to
$\mathcal{I}_t$ at every evaluated time step: neither power inequality excludes
a candidate activation vector. The feasibility layer remains part of the method,
but the active experimental restrictions are the optional-channel cardinality
constraint and the minimum dwell-time constraint.

Let $r_t$ be the remaining required activation time immediately before the action
at time $t$. If $r_t>0$, the environment executes $a_t=a_{t-1}$ regardless of
the proposed $\bar a_t$. After execution, the update is
\begin{equation}
  r_{t+1}=
  \begin{cases}
    D_{\min}-1, & a_t\neq a_{t-1},\\
    \max(r_t-1,0), & a_t=a_{t-1}.
  \end{cases}
  \label{eq:activation_timer_update}
\end{equation}
The first execution time step is therefore included in the $D_{\min}$-step
minimum. The reported values are $D_{\min}=6$ and a maximum channel warm-up of
two time steps. A newly activated channel cannot be deactivated before warm-up
completes, so $W_t=0$ structurally in this experiment.

\subsection{Partial observations and scheduler input}
\label{sec:partial_observation_update}

The environment maintains sample-and-hold values without a model-based Kalman
state estimate. Let $y_{t,i,j}$ denote a valid measurement of variable $j$
provided by ready channel $i$ at time step $t$, and let $R_{i,j}$ denote the
configured measurement-noise variance. Define $\mathcal{O}_{t,j}$ as the set
of ready channels that provide a valid measurement of variable $j$ at time
step $t$. For a scalar variable observed by at least one ready channel, the
fused value is
\begin{equation}
  \widetilde{x}_{t,j}
  =
  \frac{\sum_{i\in\mathcal{O}_{t,j}}R_{i,j}^{-1}y_{t,i,j}}
       {\sum_{i\in\mathcal{O}_{t,j}}R_{i,j}^{-1}} .
  \label{eq:observation_fusion}
\end{equation}

Wind-direction measurements are fused using a weighted circular mean. If
$\mathcal{O}_{t,j}=\varnothing$, the environment retains the previous value,
$\widetilde{x}_{t,j}=\widetilde{x}_{t-1,j}$. An observation mask indicates
whether each variable received a valid measurement at time step $t$. The
scheduler and frozen primary forecaster receive the most recent 20 value vectors
together with their observation masks. These masks distinguish new measurements
from retained sample-and-hold values.

For scalar variables, \Cref{eq:observation_fusion} assumes conditionally
independent measurement errors with variances $R_{i,j}$. Under correlated
channel errors, the rule provides a precision-weighted approximation. Computing
the minimum-variance linear estimate would require the full measurement-error
covariance matrix. The weighted circular mean for wind direction is defined
separately from the scalar fusion rule.

At each time step, the environment first updates the channel operating states
and identifies the ready channels. Measurement attempts are then performed for
the ready channels. Valid measurements are fused into the current variable
values, and variables without a valid measurement retain their previous values.
The environment subsequently updates the observation masks and channel-level
sampling ages before constructing the scheduler input.

For each channel, the scheduler receives its off, warming, or ready state,
remaining warm-up time, and normalized sampling age. The sampling age measures
the time since the channel's last sampling opportunity. It resets when an active
channel becomes ready to sample, including cases in which the stochastic
measurement attempt produces no valid measurement. Variable-level AoI follows
a different reset condition and counts the time steps since the variable's most
recent valid observation. The scheduler input further includes the previously
executed channel subset, current power ratio, diurnal phase, and compact
recent-history features.

The main policy does not receive an estimated covariance matrix. Covariance
trace is used in \Cref{prop:non_equivalence} as an alternative scheduling
objective. A diagonal predict--update uncertainty proxy is used only by the
uncertainty-objective PPO configuration.

The scheduler receives ten additional features computed deterministically from
the most recent 16 sample-and-hold value vectors and observation masks. The
features summarize recent patterns associated with particle, flux, and thermal
conditions. They also capture temporal trends and recent observation coverage.
All ten features are derived solely from information available at decision
time. Their complete definitions and clipping rules are provided in
Supplementary Section S1.

Four simulated warning signals are also available to the scheduler at decision
time. They represent particle, flux, thermal, and aggregate warning information
and include simulated noise. Condition labels are excluded from the scheduler
input. They are used to construct supervised information during training and to
group results during evaluation. A learned gate conditioned on the aggregate
warning signal combines the outputs of two feature encoders. The fused
representation is then used to compute the candidate-action scores. The critic
is implemented with a separate network. Supplementary Sections S1 and S5
provide the complete warning-signal construction and policy equations.

\begin{algorithm}[htbp]
\caption{PD-PPO policy training}
\label{alg:pdppo_training}
\small
\begin{algorithmic}
\State Train the primary forecaster on the forecaster-training partition and
freeze its parameters.
\State Evaluate the candidate fixed channel subsets on the calibration/validation partition,
compute $s_{\mathrm{calm}}$ and the event-type-specific normalization factors
$s_c$, and define the condition-specific supervised target-action mapping.
\State Pretrain the policy for 10,000 behavior cloning steps using the
supervised target actions derived from future-condition supervision labels.
\For{each PPO rollout of $n_{\mathrm{steps}}$ policy-training time steps}
  \State Initialize or continue a rollout on the policy-training partition.
  \For{each time step $t$ in the rollout}
    \State Construct $o_t$ from the observation history, channel-level runtime
    features, the previous executed activation vector, budget/time features, and
    simulated warning signals.
    \State Apply the feasibility check to obtain $\mathcal{I}_t$.
    \State Sample
    $u_t\sim\pi_{\theta}(\cdot\mid o_t,\mathcal{I}_t)$ and set
    $\bar a_t=a^{(u_t)}$.
    \State Apply the minimum dwell-time constraint to obtain $a_t$.
    \State Advance the sensing environment with $a_t$ and update the observation history.
    \State Evaluate the updated history with the frozen primary forecaster and compute
    $R_t$ using \Cref{eq:training_reward}.
    \State Store the sampled proposed-action index $u_t$, its log probability,
    the reward produced after executing $a_t$, and the training-only supervision.
  \EndFor
  \State Update the policy and critic for $n_{\mathrm{epochs}}$ passes over the
  rollout using the combined objective in \Cref{eq:pdppo_loss}.
\EndFor
\State Save the final policy checkpoint after the prescribed PPO updates and
freeze it for test evaluation.
\end{algorithmic}
\end{algorithm}

The implementation samples the policy at every training time step, including
when $r_t>0$. If the lock overrides a proposal, the rollout buffer retains the
proposed index $u_t$ and $\log\pi_{\theta_{\mathrm{old}}}
(u_t\mid o_t,\mathcal{I}_t)$, whereas the environment transition, reward, and
generalized advantage estimate are generated after executing
$a_t=a_{t-1}$. The reported training protocol therefore uses proposed-action
credit during locked time steps rather than an executed-action PPO likelihood.
The logs do not retain the counterfactual proposals needed to estimate the
override rate. This convention does not change the test losses of the frozen
final policies, but those losses do not evaluate an executed-action PPO update.

\subsection{Frozen primary forecaster and policy reward}

The primary forecaster is trained before policy optimization. Its parameters are
then frozen during policy training and evaluation.
It receives the 20-step sample-and-hold value history and its observation masks.
The time since each channel's last sampling opportunity, runtime, warning
signals, and the recent-history features above are provided only to the
scheduler. During policy training, the frozen primary forecaster returns the multi-step
forecast loss for the observation sequence produced by the current schedule.
That forecast loss forms the training loss in
\Cref{eq:subtype_normalized_training_loss}; the training loss and operational
penalties form the policy reward in \Cref{eq:training_reward}. During test
evaluation, the same frozen primary forecaster evaluates the observations produced
by the PD-PPO scheduler and every baseline.

Supplementary Table S10 separates quantities available to the online scheduler
from condition labels and future values used as training-only or evaluation-only
information. This distinction separates the privileged true-label baseline from the
PD-PPO scheduler.

The policy reward is the negative of the training loss plus the two operational
penalties used in the reported experiment.
\begin{equation}
  R_t
  =
  -\left[
  L_{\mathrm{train}}(t)
  + \lambda_{\mathrm{sw}} S_t
  + \lambda_{\mathrm{warm}} W_t
  \right],
  \label{eq:training_reward}
\end{equation}
where $L_{\mathrm{train}}$ is the condition-weighted training loss normalized
using calibration/validation quantities, $S_t$ is the normalized switching cost,
and $W_t$ is the number
of warm-up interruptions introduced at time step $t$. With $m$ logical channels,
the executed-action switching cost is
\begin{equation}
  S_t=\frac{1}{m}\lVert a_t-a_{t-1}\rVert_1.
  \label{eq:normalized_switching_cost}
\end{equation}
This normalized Hamming distance is the fraction of executed activation bits that
change between consecutive time steps. It is not a count of channel-switch
events.
Power
and startup violations are excluded by $\mathcal{I}_t$, and the minimum
dwell-time constraint determines $a_t$; these constraints therefore do
not enter the reward as residual violation terms.
The reported evaluation separates mean forecast loss from the macro-averaged
normalized forecast loss in
\Cref{eq:step_forecast_loss,eq:macro_loss}.

The training loss uses a condition-specific training weight and a normalization
factor computed on the calibration/validation partition while retaining calm
time steps and all three event types:
\begin{equation}
  L_{\mathrm{train}}(t)
  =
  \omega_{c_t}
  \frac{L_{\mathrm{step}}(t)}{s_{c_t}},
  \label{eq:subtype_normalized_training_loss}
\end{equation}
where $L_{\mathrm{step}}(t)$ is the clipped, target-weighted, scale-normalized
multi-step forecast loss returned by the frozen primary forecaster for the
current observation history, and $\omega_c\geq0$ is a condition-specific
training weight. For each event type
$c\in\{\mathrm{particle},\mathrm{flux},\mathrm{thermal}\}$, $s_c>0$ is the
median event-type-specific forecast loss across candidate fixed channel subsets
on the calibration/validation partition. The default factor $s_{\mathrm{calm}}>0$ is the
median overall forecast loss across the same candidates on that partition; it is
used for calm time steps. Thus the denominator $s_{c_t}$ in
\Cref{eq:subtype_normalized_training_loss} equals $s_c$ when $c_t$ is event type
$c$ and $s_{\mathrm{calm}}$ when $c_t$ is calm. The particle, flux, and thermal weights are 1.55, 0.95,
and 1.65, while calm time steps use unit weight. These weights and normalization
factors are fixed before policy optimization and remain unchanged during policy
training and evaluation. The current simulator-derived operating-condition label $c_t$ is
never available to the policy or critic at decision time. On the policy-training
partition, it selects the condition-specific weight and normalization factor in
\Cref{eq:subtype_normalized_training_loss}. On the test partition, it is used only
for offline loss computation and grouped reporting. Test action selection requires
the simulated warning signals but neither $c_t$ nor a reward value.

Policy learning uses two additional supervised signals formed only during
training. Define the future-condition supervision label $c_t^{\mathrm{sup}}$ as the
first non-calm simulator-derived condition label in the window from $t$ through
$t+16$, or calm when that window contains no event. This label selects a
predefined feasible supervised target action. The meteorological core channel is paired with the
shielded temperature--humidity channel in calm periods, the laser disdrometer channel for particle
events, the FC4 snow mass-flux channel for flux events, and the infrared
snow-surface-temperature channel for thermal events. These target actions supply
behavior cloning labels, while $c_t^{\mathrm{sup}}$ supervises an auxiliary
four-class future-condition classifier. Both signals are
formed only on the policy-training partition. The policy and critic receive the
simulated warning signals during policy-training rollouts and test evaluation.
Neither $c_t$ nor $c_t^{\mathrm{sup}}$ is available to the policy or critic at
decision time. The target actions are first used for 10,000 steps of
behavior cloning pretraining.
During PPO, the advantage-weighted behavior cloning term in \Cref{eq:bc_loss}
and the auxiliary future-condition classification term in
\Cref{eq:event_aux_loss} remain active.

The primary forecaster uses a TCN-style residual temporal convolutional
architecture \citep{Bai2018}. The implementation has three levels, 64 channels
per level, kernel size 3, and dropout rate 0.05. It is
fitted on the forecaster-training partition before policy training. Input and
target scales are fitted before policy updates, and the architecture is held
fixed throughout policy training and evaluation. No forecaster weights are
updated after the forecaster-training partition ends.

The forecaster is trained on observation sequences produced by full observation,
fixed channel subsets, and dynamic baselines in the forecaster-training
partition. This mixture exposes the model to dense histories and partial
histories produced by different schedules before its parameters are fixed. It
contains no calibration/validation or test observations.
Rollout counts, reset conditions, and the complete
training protocol are reported in Supplementary Section S2.

\subsection{PPO objective}

Let $o_t$ denote the scheduler input, $u_t$ the selected action index,
$\bar a_t=a^{(u_t)}$ the proposed activation vector, $a_t$ the executed
activation vector, and $\hat{A}_t$ the advantage estimate. The policy defines a
categorical distribution over the feasible action set,
\begin{equation}
  \pi_{\theta}(u\mid o_t,\mathcal{I}_t)
  =
  \frac{\exp z_{\theta}(o_t,u)}
  {\sum_{v\in\mathcal{I}_t}\exp z_{\theta}(o_t,v)}
  \mathbf{1}\{u\in\mathcal{I}_t\},
  \label{eq:masked_policy}
\end{equation}
where $z_{\theta}(o_t,u)$ is the policy score for activation vector $a^{(u)}$ and
$\mathcal{I}_t$ is the power- and startup-feasible index set. The activation
vector containing only the mandatory meteorological core channel belongs to the
candidate action set and satisfies both limits, so
$\mathcal{I}_t$ is nonempty at every time step.
PPO updates this distribution with the clipped surrogate
\begin{equation}
  \mathcal{L}_{\mathrm{clip}}(\theta)
  =
  -\mathbb{E}_t
  \left[
  \min\left(
  \rho_t(\theta)\hat{A}_t,\,
  \operatorname{clip}\!\left(\rho_t(\theta),1-\epsilon,1+\epsilon\right)
  \hat{A}_t
  \right)
  \right],
  \label{eq:ppo_clip}
\end{equation}
where
\begin{equation}
  \rho_t(\theta)
  =
  \frac{\pi_{\theta}(u_t\mid o_t,\mathcal{I}_t)}
       {\pi_{\theta_{\mathrm{old}}}(u_t\mid o_t,\mathcal{I}_t)} .
\end{equation}
The reported PPO loss contains the PPO surrogate, value loss, entropy
regularization, and two auxiliary terms tied to the forecast training setup.
\begin{equation}
  \mathcal{L}_{\mathrm{PPO}}
  =
  \mathcal{L}_{\mathrm{clip}}
  + c_v\mathcal{L}_{V}
  - c_H\mathcal{H}\!\left(\pi_{\theta}(\cdot\mid o_t,\mathcal{I}_t)\right)
  + \beta_{\mathrm{BC}}\mathcal{L}_{\mathrm{BC}}
  + \beta_{\mathrm{cond}}\mathcal{L}_{\mathrm{cond}} .
  \label{eq:pdppo_loss}
\end{equation}
For labeled training examples, let $u_t^{\mathrm{target}}\in\mathcal{I}_t$
denote the feasible supervised target action described above. The
advantage-weighted behavior cloning term is
\begin{equation}
  \mathcal{L}_{\mathrm{BC}}
  =
  -
  \frac{
  \sum_t \ell_t[\hat{A}_t]_+\,\log
  \pi_{\theta}(u_t^{\mathrm{target}}\mid o_t,\mathcal{I}_t)
  }{
  \sum_t \ell_t[\hat{A}_t]_+ + \varepsilon
  },
  \label{eq:bc_loss}
\end{equation}
where $[x]_+=\max(x,0)$ and $\ell_t\in\{0,1\}$ marks examples with a valid
supervised target. Let $\eta_t\in\{0,1\}$ mark examples for which the
future-condition supervision label is available to the auxiliary classifier. An
auxiliary future-condition classification loss is applied to the auxiliary
future-condition classifier. Let
$h_t$ be the policy representation, $q_{\theta}(c\mid h_t)$ the auxiliary
probability of operating condition
$c\in\{\mathrm{calm},\mathrm{particle},\mathrm{flux},\mathrm{thermal}\}$, and
$c_t^{\mathrm{sup}}$ the future-condition supervision label used only during policy
training. The reported stride-one configuration constructs both supervised
targets from the same $c_t^{\mathrm{sup}}$ and sets
$\eta_t=\ell_t=1$ at every PPO step. The general
implementation constructs the indicators separately. The auxiliary future-condition classification loss is
\begin{equation}
  \mathcal{L}_{\mathrm{cond}}
  =
  -
  \frac{
  \sum_t \eta_t \log q_{\theta}(c_t^{\mathrm{sup}}\mid h_t)
  }{
  \sum_t \eta_t + \varepsilon
  } .
  \label{eq:event_aux_loss}
\end{equation}

Supplementary Table S1 lists the full protocol and training hyperparameters.
\subsection{Reward-function and algorithm comparisons}
\label{sec:matched_controls}

The AoI-objective and uncertainty-objective PPO configurations retain the same PPO
policy and critic architecture, scheduler input, candidate action set, feasibility check,
chronological partitions, condition-specific weights, supervised target actions
available only during training, and auxiliary future-condition classification loss, but
replace the scalar forecast-loss objective. The AoI-objective PPO variant averages normalized
observed-variable age over the observation history; the uncertainty-objective
configuration averages a bounded diagonal predict--update variance proxy. Their
recursions and bounds are given in
Supplementary Section S3.
Each scalar objective supplies the corresponding rollout reward, from which the
value targets and advantages are computed. Replacing it therefore changes the
PPO surrogate, critic target, and the positive-advantage weights in
$\mathcal{L}_{\mathrm{BC}}$. These are matched PPO training configurations with
different scalar objectives, not isolated scalar-reward ablations.

The Double DQN training configuration uses a dueling network and the same
feasible-action mask, state representation, reward, action set, constraints,
policy-training partition, and test windows as the PD-PPO scheduler. Online and
target networks implement the
Double DQN target over feasible next actions, with three-step returns and
experience replay. This configuration does not use behavior cloning pretraining, the continuing
advantage-weighted behavior cloning term, or the auxiliary future-condition
classifier. Its complete settings are listed in
Supplementary Table S1.

\subsection{Chronological data split and baselines}

Each simulated ground-truth time series is split into forecaster-training,
policy-training, calibration/validation, and test partitions. The
calibration/validation partition computes $s_c$ and $s_{\mathrm{calm}}$, defines
the condition-to-action mapping used to construct supervised targets on
policy-training data, and selects the fixed-schedule baseline. Test windows are
not used for forecaster training, policy updates,
supervised-target generation, fixed-baseline selection, or normalization-factor
fitting.

The calibration/validation interval follows the policy-training interval
chronologically. Computationally, however, its normalization factors,
condition-to-action mapping, and fixed-schedule selection are computed offline
before policy optimization and then held fixed. Calibration/validation
transitions are not used as policy-training rollouts.

The chronological split is summarized in Supplementary Fig. S1.

The fixed-schedule baseline tests whether adaptation is needed. The AoI-priority
baseline, round-robin baseline, and random-feasible baseline provide non-learned
dynamic comparisons under the same operating rules. The privileged one-step
look-ahead reference using next-step test targets chooses the currently best feasible action;
those targets are unavailable to the deployed scheduler. Its one-step objective omits
minimum dwell-time consequences, warm-up, and observation age. It therefore does not isolate the
value of sequential decision making. The warning-rule baseline maps simulated
warning signals to channel subsets selected on the calibration/validation
partition. The privileged true-label baseline receives the exact sequence
$c_t,\ldots,c_{t+16}$ at each decision; these condition labels are unavailable
to the deployed scheduler. Neither privileged baseline is a guaranteed upper
bound.

Supplementary Table S11 fixes the role of each baseline before the results
section.

Exact decision rules, calibration/validation choices, tie breaks, and information availability
for these baselines are listed in Supplementary Section S6.
